% Corrigé intégré automatiquement pour la banque Exercices.
% Source : A_Integrer/DDS_04/066_Modelisation_Geometrie/066_Modelisation_Geometrie.tex
\begin{corrigebox}[Corrigé]
\CorrigeQuestion{1}
La longueur initiale lambda zéro est la racine carrée de la somme de L au carré et de R au carré. La longueur lambda est la racine carrée de la somme de L au carré et de la différence entre R et x au carré.
\CorrigeQuestion{2}
Deux fois la différence entre lambda zéro et lambda est égale au produit de R par la variation angulaire theta.
\CorrigeQuestion{3}
Le déplacement x est égal à R diminué de la racine carrée de l'expression obtenue à partir de R au carré, du produit de R par la longueur initiale et par la variation angulaire, puis du terme quadratique associé à cette variation.
\CorrigeQuestion{4}
On isole le solide ou l'ensemble indiqué, on dresse le bilan complet des actions mécaniques extérieures et on choisit un point de réduction qui élimine le maximum d'inconnues. Le principe fondamental de la statique fournit ensuite les équations de résultante et de moment ; les projections utiles donnent les inconnues, puis leur signe permet de vérifier le sens supposé.
\CorrigeQuestion{5}
Le coefficient K est égal à la moitié de la racine carrée de la somme de L au carré et de R au carré.
\end{corrigebox}
